\documentclass[conference]{IEEEtran}
\usepackage{cite}
\usepackage{graphicx}
\usepackage{hyperref}
\usepackage{amsmath}
\usepackage{amssymb}

\title{Generative Models for Rare Disease Biomedical Data Synthesis: A Literature Review}
\author{
    % Add author information here
}

\begin{document}
\maketitle

\begin{abstract}
Rare diseases affect millions globally but suffer from extreme data scarcity, creating significant barriers to developing effective diagnostic and therapeutic approaches. This review examines how generative models—including Generative Adversarial Networks (GANs), Variational Autoencoders (VAEs), and diffusion models—address the challenge of synthesizing biomedical data for ultra-rare conditions. We analyze technical approaches, evaluate effectiveness in mitigating data scarcity and class imbalance, and examine applications across electronic health records, multi-omics profiles, and medical imaging modalities.
\end{abstract}

\section{Introduction}

Rare diseases affect approximately 350 million individuals globally across nearly 7,000 distinct conditions, yet each condition represents less than one percent of the overall disease burden \cite{finetti2025data}. This epidemiological paradox creates substantial barriers to clinical research and diagnostic development. Patients experience an average diagnostic delay of six years from symptom onset to accurate diagnosis, during which time they navigate fragmented healthcare systems with limited access to specialized expertise. The combination of small patient cohorts, geographic dispersion, and phenotypic heterogeneity results in datasets ranging from dozens to a few thousand cases---far below the sample sizes required for robust machine learning model training \cite{finetti2025data,mendes2025synthetic}.

Traditional approaches to data augmentation, including geometric transformations and photometric adjustments, have dominated biomedical applications due to their interpretability and computational efficiency. A recent scoping review of 118 studies published between 2018 and 2025 found that classical augmentation techniques accounted for 49.3\% of data expansion strategies in rare disease research \cite{finetti2025data}. However, these methods are fundamentally constrained by the variability present in the original dataset, limiting their capacity to generate novel yet biologically plausible samples. As artificial intelligence applications in healthcare have expanded, the need for more sophisticated data generation approaches has become increasingly apparent.

Deep generative models have emerged as a promising solution to overcome the limitations of classical augmentation. Generative Adversarial Networks (GANs), which employ adversarial training between generator and discriminator networks, have demonstrated particular success in medical imaging applications. Frid-Adar et al. achieved a 7.1 percentage point improvement in liver lesion classification sensitivity through GAN-based synthetic CT image generation, with final performance reaching 85.7\% sensitivity and 92.4\% specificity \cite{fridadar2018gan}. The application of GANs has extended beyond imaging to electronic health records (EHRs), where medGAN successfully synthesized multi-label discrete patient records with 97.77\% accuracy while preserving privacy through fully synthetic data generation \cite{choi2017generating}.

Variational Autoencoders (VAEs) offer an alternative approach through probabilistic modeling of latent distributions, demonstrating particular utility in scenarios with limited computational resources. While VAEs may produce slightly less sharp images compared to GANs, their lower computational requirements and reduced risk of mode collapse make them valuable for resource-constrained settings \cite{mendes2025synthetic}. More recently, diffusion-based models have introduced novel capabilities for sequential data generation. The MedDiffusion framework, applying denoising diffusion probabilistic models to EHR data for the first time, achieved 77.88\% PR-AUC on kidney disease risk prediction by generating synthetic patient visits conditioned on medical history \cite{zhong2024meddiffusion}.

Despite these advances, substantial challenges remain. Finetti et al. identified that over 40\% of published studies lack sufficient quantitative detail regarding dataset sizes before and after augmentation, hindering reproducibility and comparative analysis \cite{finetti2025data}. Privacy concerns persist even with synthetic data generation, as membership inference and model inversion attacks remain potential vulnerabilities \cite{mendes2025synthetic}. Validation practices vary widely, with most studies evaluating synthetic data quality solely through downstream model performance rather than independent biological plausibility assessment. The computational demands of deep generative models, requiring GPU clusters and extensive hyperparameter tuning, create additional barriers to widespread adoption.

This review synthesizes current evidence on generative model applications for rare disease biomedical data synthesis. We examine the technical foundations of GANs, VAEs, and diffusion models in Section~\ref{sec:methods}, analyze strategies for mitigating data scarcity and class imbalance in Section~\ref{sec:scarcity}, review modality-specific applications across imaging, EHRs, and omics data in Section~\ref{sec:applications}, discuss current limitations in Section~\ref{sec:limitations}, and conclude with future directions for clinical integration in Section~\ref{sec:conclusion}. Our focus remains on how these generative approaches specifically address the unique constraints imposed by rare disease research, where traditional data collection strategies fail to accumulate sufficient samples for modern machine learning pipelines.


\section{Generative Models in Biomedical Contexts}
\label{sec:methods}

The application of deep generative models to biomedical data synthesis requires careful consideration of architectural choices, training dynamics, and domain-specific constraints. This section examines the technical foundations of three predominant approaches---Generative Adversarial Networks (GANs), Variational Autoencoders (VAEs), and diffusion models---with emphasis on adaptations that address the unique characteristics of rare disease datasets.

\subsection{Generative Adversarial Networks}

GANs employ adversarial training between two neural networks: a generator $G$ that produces synthetic samples from random noise, and a discriminator $D$ that distinguishes real from generated data. The training objective formulates a minimax game where the generator seeks to minimize the discriminator's classification accuracy while the discriminator maximizes it. This adversarial process, when successful, drives the generator to produce samples indistinguishable from the real data distribution.

Frid-Adar et al. demonstrated the application of Deep Convolutional GANs (DCGANs) to medical imaging, implementing a generator architecture with four fractionally-strided convolutional layers using 5×5 kernels \cite{fridadar2018gan}. The network transforms 100-dimensional uniform noise vectors into 64×64 pixel liver lesion images through progressive upsampling. Batch normalization and ReLU activations stabilize training, while the final layer employs tanh activation to constrain output values to the appropriate intensity range. The discriminator mirrors this architecture in reverse, classifying input images as real or synthetic through successive downsampling operations.

A critical challenge in GAN training is mode collapse, where the generator produces limited diversity by concentrating on a subset of the target distribution. Choi et al. introduced minibatch averaging as an effective mitigation strategy for electronic health record synthesis \cite{choi2017generating}. Rather than processing individual samples independently, the discriminator computes feature statistics across minibatches and incorporates this aggregated information into its classification decision. This technique encourages the generator to produce diverse outputs without the hyperparameter sensitivity of alternative approaches such as minibatch discrimination. The medGAN architecture further incorporates shortcut connections between generator layers, enabling more effective gradient flow during backpropagation and accelerating convergence.

The choice between different GAN variants depends on the synthesis task requirements. Auxiliary Classifier GANs (ACGANs) extend the discriminator to predict class labels in addition to authenticity, enabling conditional generation of specific disease subtypes. Frid-Adar et al. found that DCGAN outperformed ACGAN for their liver lesion classification task, suggesting that the additional auxiliary loss may not always improve performance when category information is already well-represented in the training data \cite{fridadar2018gan}. Conditional GANs (cGANs) and their variants---including CycleGANs for unpaired image translation and StyleGANs for fine-grained control over generated features---have each found distinct applications in biomedical synthesis \cite{mendes2025synthetic}.

\subsection{Variational Autoencoders}

VAEs approach generative modeling through probabilistic inference, learning a latent representation that captures the underlying data distribution. The architecture comprises an encoder network that maps input data to a probability distribution over a lower-dimensional latent space, and a decoder that reconstructs samples from latent vectors. Unlike GANs' adversarial objective, VAEs optimize a variational lower bound on the data log-likelihood, combining a reconstruction loss with a Kullback-Leibler divergence term that regularizes the latent distribution toward a standard normal prior.

This probabilistic formulation provides theoretical guarantees on the quality of the learned latent space and eliminates the training instability often encountered with GANs. The regularization term ensures that similar data points map to nearby latent representations, facilitating interpolation between samples and enabling controlled generation through latent space navigation. Conditional VAEs (CVAEs) extend this framework by conditioning both encoder and decoder on auxiliary information such as disease labels, demonstrating particular effectiveness with smaller datasets that require class-specific generation \cite{mendes2025synthetic}.

The primary limitation of VAEs relative to GANs lies in reconstruction quality. The continuous latent space and probabilistic decoder produce samples that, while statistically faithful to the training distribution, often exhibit reduced sharpness compared to adversarially-trained generators. For medical imaging applications requiring high spatial resolution, this trade-off manifests as slightly blurred synthetic images. However, VAEs require substantially lower computational resources than GANs, making them viable for resource-constrained environments \cite{mendes2025synthetic}. Hybrid VAE-GAN architectures attempt to combine the advantages of both approaches, using VAE structure for stable training while incorporating adversarial objectives to sharpen generated outputs.

\subsection{Diffusion Models}

Diffusion probabilistic models represent an emerging paradigm for generative modeling, formulating synthesis as a progressive denoising process. Training proceeds in two phases: a forward diffusion process gradually adds Gaussian noise to real data over $T$ timesteps until reaching pure noise, followed by learning a reverse process that reconstructs clean data from noisy inputs. The reverse process, parameterized by a neural network, predicts the noise component at each timestep, enabling iterative refinement from random samples to high-quality outputs.

Zhong et al. adapted this framework to sequential electronic health record data through the MedDiffusion architecture \cite{zhong2024meddiffusion}. The model embeds medical diagnosis codes and temporal information into visit representations using a combination of learned embeddings and gap-based time encodings. A bidirectional Long Short-Term Memory (LSTM) network processes these visit sequences to capture longitudinal patient trajectories. The key innovation lies in the step-wise attention mechanism, which aggregates information from the current visit and previous hidden states to condition the diffusion process on patient history. This conditioning enables generation of synthetic visits that maintain realistic disease progression patterns while expanding the training dataset.

The forward diffusion process in MedDiffusion adds noise to visit embeddings according to a predetermined variance schedule, gradually corrupting the structured medical code representations. The reverse process trains a denoising network to predict the original clean embedding from noisy inputs at each timestep, conditioned on the LSTM-encoded patient history. Generating a new synthetic visit begins with Gaussian noise and applies the learned denoising steps iteratively, with the LSTM context ensuring temporal consistency across the patient trajectory.

Zhong et al. demonstrated that this diffusion-based approach outperformed GAN baselines (ehrGAN, TabDDPM) across multiple health risk prediction tasks, achieving 77.88\% PR-AUC on kidney disease prediction compared to 74.92\% for the best GAN baseline \cite{zhong2024meddiffusion}. The superior performance stems from diffusion models' ability to generate in continuous latent space while avoiding the mode collapse and training instability that plague GAN training on sequential discrete data. Ablation studies confirmed that both the step-wise attention mechanism and the diffusion-based generation contribute essential components to the overall performance.

\subsection{Comparative Analysis}

The choice of generative model architecture depends on multiple factors including data modality, available computational resources, dataset size, and synthesis objectives. Medical imaging applications, which dominated 70.3\% of reviewed studies, predominantly employ GANs due to their capacity for generating sharp, high-resolution images \cite{finetti2025data}. The adversarial training objective proves particularly effective for capturing the complex spatial dependencies present in CT scans, MRI sequences, and histopathology slides. However, this comes at substantial computational cost---GAN training typically requires GPU clusters and careful hyperparameter tuning to achieve stable convergence.

VAEs offer a compelling alternative when researchers face computational resource constraints or require theoretical guarantees on the latent space structure. Their lower training complexity and reduced risk of mode collapse make them particularly suitable for exploratory analysis and scenarios where slightly reduced image sharpness is acceptable relative to training efficiency. The probabilistic framework also facilitates uncertainty quantification, enabling assessment of the model's confidence in generated samples.

Diffusion models have emerged most recently and show particular promise for sequential and temporal data where maintaining consistency across time steps is critical. The iterative denoising process naturally incorporates contextual information, making diffusion models well-suited to electronic health record synthesis and longitudinal patient trajectory generation. While diffusion models require more inference steps than single-pass GAN or VAE generation, recent advances in accelerated sampling have reduced this computational overhead substantially \cite{zhong2024meddiffusion}.

Hybrid approaches combining elements from multiple frameworks represent an active area of investigation. VAE-GAN architectures leverage VAE encoders for stable latent space learning while incorporating GAN discriminators to sharpen outputs. Conditional variants of all three model classes enable class-specific generation, proving essential for rare disease applications where maintaining disease subtype distinctions is critical. The continued rapid development of generative modeling techniques suggests that optimal architectures for biomedical synthesis remain an evolving target rather than a settled question.

\section{Mitigating Scarcity and Imbalance}
\label{sec:scarcity}

Data scarcity in rare disease research manifests in multiple dimensions beyond simple sample counts. Small cohorts exhibit limited phenotypic variability, constraining model exposure to the full spectrum of disease presentations. Class imbalance, where rare disease cases constitute a small fraction relative to healthy controls or common conditions, further compounds training challenges. This section examines strategies through which generative models address these interconnected challenges, with emphasis on quantitative expansion rates, quality assessment methods, and approaches to preserving biological plausibility.

\subsection{Dataset Expansion Strategies}

The most direct application of generative models addresses absolute sample scarcity through synthetic data creation. Finetti et al. analyzed expansion rates across 118 studies, finding that medical imaging datasets achieved an average threefold increase through classical and deep generative augmentation \cite{finetti2025data}. However, this aggregate metric masks substantial variability across applications and techniques. Frid-Adar et al. expanded an initial dataset of 182 liver lesion CT scans through classical augmentation (rotations, translations, flips, scaling) to approximately 30,000 augmented samples per lesion class before applying GAN synthesis \cite{fridadar2018gan}. The combination of classical and deep generative approaches enabled training convolutional neural networks that would otherwise fail to converge on the original small dataset.

The choice of expansion magnitude depends on balancing dataset diversity against computational feasibility and the risk of introducing artifacts. Excessive reliance on synthetic data can lead to models that perform well on generated samples but fail to generalize to real-world variation. Mendes et al. documented synthetic cohort expansions ranging from threefold to tenfold across multiple rare disease applications \cite{mendes2025synthetic}. For acute myeloid leukemia studies, CTAB-GAN+ generated synthetic patient records that tripled the available training data while maintaining the statistical characteristics of rare genetic variants. The key distinction from classical augmentation lies in the generative models' capacity to produce novel combinations of features rather than simple transformations of existing samples.

Temporal and sequential data present unique expansion challenges due to the need for maintaining realistic longitudinal patterns. Zhong et al. addressed this through conditional generation in MedDiffusion, where the model generates synthetic patient visits based on medical history rather than independently \cite{zhong2024meddiffusion}. This approach ensures that synthetic disease trajectories exhibit plausible progression patterns consistent with clinical understanding. The model generated synthetic visits for patients with kidney disease, chronic obstructive pulmonary disease (COPD), and amnesia, expanding datasets with 2,810 to 7,314 positive cases by factors ranging from two to five depending on the specific prediction task requirements.

\subsection{Quality Assessment and Validation}

A critical challenge in synthetic data generation is establishing appropriate quality metrics beyond downstream model performance. Finetti et al. identified that over 40\% of reviewed studies provided insufficient quantitative detail regarding dataset sizes and expansion rates, while fewer than 33\% reported both pre- and post-augmentation metrics \cite{finetti2025data}. This lack of standardized reporting complicates comparative analysis and reproducibility assessment across studies.

Visual assessment by domain experts has emerged as a common validation approach for medical imaging applications. Frid-Adar et al. had radiologists evaluate synthetic liver lesions for anatomical plausibility and diagnostic utility \cite{fridadar2018gan}. The experts confirmed that GAN-generated lesions exhibited realistic texture patterns and morphological characteristics across all three lesion classes (cysts, metastases, hemangiomas). However, this validation method scales poorly and introduces subjectivity, as inter-rater reliability may vary and systematic evaluation protocols remain underdeveloped.

Statistical validation through distribution matching provides more objective assessment. Choi et al. demonstrated that medGAN-generated electronic health records preserved the marginal distributions of diagnosis and medication codes while maintaining appropriate correlations between co-occurring conditions \cite{choi2017generating}. The synthetic data reproduced the prevalence of rare conditions within the expected statistical variation of the real dataset. Tests comparing mean, variance, and higher-order moments across real and synthetic distributions confirmed statistical fidelity. Privacy assessment through membership inference attacks and attribute disclosure metrics further validated that synthetic records did not memorize individual patients from the training data.

Performance-based validation evaluates synthetic data utility through downstream prediction tasks. Frid-Adar et al. measured liver lesion classification performance with and without synthetic augmentation, observing improvements from 78.6\% to 85.7\% sensitivity and from 88.4\% to 92.4\% specificity \cite{fridadar2018gan}. These 7.1 and 4.0 percentage point gains demonstrate that synthetic data meaningfully enhanced model generalization. Similarly, Mendes et al. reported that combining synthetic and real brain MRI data achieved 85.9\% classification accuracy, with Dice score improvements ranging from 3\% to 15\% for segmentation robustness \cite{mendes2025synthetic}. Zhong et al. showed that MedDiffusion improved health risk prediction PR-AUC from 71.05\% to 77.88\% for kidney disease, representing a 6.83 percentage point gain over the best non-augmented baseline \cite{zhong2024meddiffusion}.

\subsection{Addressing Class Imbalance}

Class imbalance, where rare disease cases are substantially outnumbered by controls, presents distinct challenges from overall data scarcity. Traditional oversampling techniques such as SMOTE (Synthetic Minority Over-sampling Technique) generate synthetic minority class examples through interpolation between existing samples. While effective for tabular data with limited dimensionality, SMOTE struggles with high-dimensional medical imaging and complex sequential records. Deep generative models offer more sophisticated alternatives by learning the underlying distribution of the minority class and generating novel examples that extend beyond linear interpolation.

Conditional generation mechanisms enable targeted synthesis of underrepresented disease subtypes. Conditional GANs and VAEs accept class labels as additional inputs, directing the generator to produce samples from specific diagnostic categories. This proves particularly valuable in rare disease contexts where certain phenotypic variants or genetic subtypes may have fewer than ten documented cases. Finetti et al. documented applications in hematological malignancies, where acute lymphoblastic leukemia subtypes exhibit highly imbalanced class distributions \cite{finetti2025data}. Synthetic generation of underrepresented subtypes enabled training classifiers that would otherwise collapse to predicting only the majority class.

The effectiveness of class balancing through synthetic generation depends critically on the original training set containing sufficient minority class examples to capture distributional characteristics. When fewer than 20-30 examples exist, even sophisticated generative models risk producing samples that merely reproduce training set variation rather than capturing the true population distribution. In such extreme scarcity scenarios, transfer learning from related conditions or incorporating prior knowledge through rule-based constraints may provide necessary inductive biases. Mendes et al. described approaches combining synthetic generation with domain expert input to ensure generated samples respect known biological constraints \cite{mendes2025synthetic}.

\subsection{Biological Plausibility and Artifact Detection}

A persistent concern with synthetic biomedical data is the potential generation of biologically implausible samples that, while statistically consistent with training data, violate known physiological or anatomical principles. Generative models trained purely on data-driven objectives lack explicit mechanisms to enforce domain constraints. A GAN might generate a liver lesion with impossible tissue density values or an electronic health record with contradictory diagnosis codes that never co-occur in clinical practice.

Several approaches attempt to address this limitation. Hybrid methods combine data-driven generation with rule-based systems that filter or constrain outputs according to known medical knowledge. For example, generated electronic health records can be validated against clinical guidelines that specify valid medication combinations and contraindicated prescription patterns. Mendes et al. emphasized the importance of incorporating domain expertise throughout the generation pipeline rather than treating validation as a post-hoc step \cite{mendes2025synthetic}.

Adversarial validation provides another quality assurance mechanism, where a separate classifier attempts to distinguish synthetic from real samples. If the discriminator achieves high accuracy, this indicates that synthetic samples exhibit detectable artifacts or systematic differences from real data. However, this approach fundamentally measures statistical similarity rather than biological plausibility---a model might generate statistically indistinguishable samples that nonetheless violate medical constraints not represented in the training distribution.

The computational cost of synthetic data generation also warrants consideration in resource allocation decisions. Finetti et al. noted that deep generative model training typically requires GPU clusters and substantial hyperparameter tuning effort \cite{finetti2025data}. For datasets with fewer than 100 samples, the overhead of GAN training may exceed the benefit relative to simpler augmentation techniques. The decision between classical augmentation, VAEs, GANs, or diffusion models must weigh not only performance metrics but also practical constraints on computational resources, domain expertise availability, and time-to-deployment requirements.

\section{Modality-Specific Applications}
\label{sec:applications}

Generative model architectures require adaptation to the distinct characteristics of different biomedical data modalities. Medical imaging, electronic health records, and multi-omics data each present unique challenges in terms of dimensionality, structure, temporal dependencies, and domain constraints. This section examines how researchers have tailored generative approaches to specific data types, with particular attention to the interplay between modality characteristics and model design choices.

\subsection{Medical Imaging}

Medical imaging applications have dominated the deployment of generative models in rare disease research, accounting for 70.3\% of reviewed studies \cite{finetti2025data}. This concentration reflects both the availability of established deep learning architectures for image data and the acute need for expanded training sets in radiology and pathology where acquisition costs and patient scarcity severely limit dataset sizes.

Frid-Adar et al. applied Deep Convolutional GANs to computed tomography (CT) imaging for liver lesion classification \cite{fridadar2018gan}. Starting from 182 lesion samples across three diagnostic categories (cysts, metastases, hemangiomas), they generated synthetic 64×64 pixel regions of interest that captured characteristic texture patterns and morphological features. The generator architecture employed fractionally-strided convolutions to progressively upsample from latent noise vectors to full-resolution images, with batch normalization stabilizing training dynamics. Classical augmentation techniques (rotations, translations, scaling, flipping) expanded each original lesion to approximately 30,000 transformed samples before GAN synthesis added further diversity. The combined augmentation strategy enabled convolutional neural network training that achieved 85.7\% sensitivity and 92.4\% specificity, representing 7.1 and 4.0 percentage point improvements over classical augmentation alone.

Expert radiologist validation confirmed the diagnostic utility of synthetic lesions, though systematic evaluation protocols remain underdeveloped. Visual inspection verified that generated samples exhibited realistic tissue density values, appropriate spatial relationships between anatomical structures, and diagnostically relevant features. However, this validation approach scales poorly to large synthetic datasets and introduces subjectivity that complicates reproducibility assessment across institutions.

The application of generative models to medical imaging extends beyond CT to magnetic resonance imaging (MRI), retinal fundus photography, and histopathology. Mendes et al. documented that combining synthetic and real brain MRI data achieved 85.9\% classification accuracy with Dice score improvements of 3-15\% for segmentation tasks \cite{mendes2025synthetic}. For age-related macular degeneration, ResNet-18 architectures trained on GAN-augmented retinal images reached 85\% diagnostic accuracy, outperforming human expert consensus at 75-80\% \cite{mendes2025synthetic}. These results demonstrate that synthetic augmentation can surpass simple dataset expansion to enable superhuman performance in specific diagnostic contexts.

The choice between GAN variants depends on the specific imaging task. Conditional GANs enable generation of specific lesion types or disease stages by incorporating class labels as additional inputs. CycleGANs facilitate unpaired image-to-image translation, useful for modality conversion or artifact removal when paired training examples are unavailable. StyleGANs provide fine-grained control over generated image attributes through manipulation of latent style vectors at different resolution scales. The continued proliferation of GAN architectures specialized for different imaging modalities reflects the domain-specific nature of optimal generative model design.

\subsection{Electronic Health Records}

Electronic health record (EHR) data represented 19.5\% of reviewed applications \cite{finetti2025data}, encompassing both cross-sectional aggregated features and longitudinal event sequences. The discrete, high-dimensional, and sparse nature of EHR data necessitates distinct architectural adaptations compared to continuous-valued imaging data.

Choi et al. developed medGAN specifically for synthesizing multi-label discrete patient records \cite{choi2017generating}. The architecture combines an autoencoder for dimensionality reduction with a GAN for generating samples in the learned latent space. The autoencoder compresses binary and count variables representing diagnosis and medication codes into continuous latent vectors, enabling smoother GAN training than operating directly on discrete features. The generator incorporates shortcut connections that propagate gradients more effectively during backpropagation, while the discriminator employs minibatch averaging to combat mode collapse without hyperparameter tuning.

Testing on three EHR datasets ranging from 30,738 to 258,559 patients demonstrated that synthetic records preserved marginal distributions and co-occurrence patterns of medical codes. Disease prediction models trained on synthetic data achieved 97.77\% accuracy, with sensitivity of 95.21\% and specificity of 99.25\%. Privacy analysis through membership inference attacks and attribute disclosure metrics confirmed that synthetic records did not memorize individual training samples, enabling data sharing between institutions without residual re-identification risk.

The temporal dimension of EHR data introduces additional complexity beyond cross-sectional aggregation. Disease progression follows characteristic trajectories that synthetic data must respect to maintain clinical validity. Zhong et al. addressed this challenge through MedDiffusion, adapting diffusion probabilistic models to sequential medical records \cite{zhong2024meddiffusion}. The architecture embeds diagnosis codes and visit timing information into continuous representations processed by bidirectional LSTMs that capture longitudinal dependencies. A step-wise attention mechanism conditions synthetic visit generation on patient history, ensuring temporal consistency in disease progression patterns.

MedDiffusion outperformed GAN baselines across four health risk prediction tasks, achieving 77.88\% PR-AUC for kidney disease prediction compared to 74.92\% for the best GAN alternative. The performance gain derives from diffusion models' superior handling of discrete sequential data and their avoidance of the mode collapse that frequently afflicts GAN training on structured medical records. The conditioning mechanism proved critical---ablation studies removing the step-wise attention or the diffusion process substantially degraded performance, confirming that both components contribute essential capabilities.

The expansion of EHR datasets through synthetic generation addresses not only quantity limitations but also representation biases. Real-world EHR systems overrepresent frequently-encountered conditions and patient demographics while underrepresenting rare diseases and minority populations. Synthetic generation with appropriate conditioning can rebalance these distributions, though careful validation remains necessary to ensure that generated samples reflect realistic disease prevalence rather than artifacts of the generation process.

\subsection{Rare Genetic Diseases and Knowledge-Guided Synthesis}

Multi-omics data and rare genetic diseases present extreme scarcity challenges where even dozens of patient samples may be unavailable. These applications constitute only 5.1\% of reviewed studies \cite{finetti2025data}, yet represent critical domains where generative approaches could provide disproportionate impact. The ultra-low sample regime necessitates integration of domain knowledge beyond purely data-driven learning.

Alsentzer et al. demonstrated an alternative paradigm through SHEPHERD, a few-shot learning approach for phenotype-driven diagnosis that relies primarily on synthetic patient data \cite{alsentzer2025few}. Rather than augmenting small real datasets, SHEPHERD trains on 42,624 simulated patients representing 2,132 unique Mendelian disorders. The synthesis process grounds generation in biomedical knowledge graphs encoding relationships between phenotypes, genes, diseases, and biological pathways. Each synthetic patient comprises phenotypic features sampled from disease-specific distributions in the Orphanet rare disease database, with additional corruption and noise injection to simulate diagnostic uncertainty and incomplete clinical information.

The SHEPHERD architecture employs graph neural networks that operate directly on patient phenotype subgraphs extracted from a heterogeneous biomedical knowledge graph containing 105,220 nodes and over 1 million edges. Graph attention networks aggregate neighborhood information through relation-specific message passing, enabling the model to leverage both direct and indirect associations between genes and phenotypes. This proves critical for rare diseases where only 28\% of patient phenotypic features have known direct associations with causal genes. For patients with phenotypes separated by more than two relationship hops from their causal gene, SHEPHERD achieved 77.8\% accuracy in ranking the correct gene within the top 5 predictions.

Validation on three independent real-world cohorts totaling 2,042 patients with 378 unique genetic conditions demonstrated that synthetic training translates to real-world diagnostic utility. SHEPHERD ranked the correct causal gene first for 40\% of Undiagnosed Diseases Network patients when provided expert-curated candidate gene lists, achieving 2.2-fold improvement over non-guided baselines. Perhaps most remarkably, 81.8\% of real patient phenotype terms appeared in the synthetic training cohort despite the entirely simulated origin, and synthetic patients clustered with real patients of the same disease category in learned embedding spaces.

This knowledge-guided synthesis paradigm differs fundamentally from pure generative modeling. Rather than learning distributions solely from data, the approach encodes structured domain knowledge that constrains generation to biologically plausible combinations. For rare genetic diseases where acquiring thousands of real patients is impossible, this knowledge integration enables deep learning models to achieve clinical utility despite extreme data scarcity. The 20 synthetic patients generated per disease in SHEPHERD's training set would be insufficient for traditional supervised learning but prove adequate when combined with knowledge graph structure and few-shot learning objectives.

Genomic data synthesis extends beyond phenotype-level modeling to include DNA/RNA sequence generation and gene expression profile synthesis. Mendes et al. described applications of GANs and VAEs to capturing rare genetic variants that appear infrequently in population databases \cite{mendes2025synthetic}. CTAB-GAN+ and normalizing flows have generated synthetic patient cohorts for acute myeloid leukemia studies, achieving threefold expansion while preserving the statistical characteristics of rare allelic combinations. These synthetic genomes enable variant calling algorithm development and population genetics analyses without requiring consent protocols or privacy protection mechanisms that constrain real genomic data sharing.

The integration of generative models with knowledge graphs, ontologies, and mechanistic biological models represents a frontier for rare disease applications. Pure data-driven approaches struggle when training examples number in the single or low double digits. Hybrid methods that incorporate known gene-disease relationships, protein-protein interaction networks, metabolic pathway structures, and phenotype ontologies can generate synthetic data that respects biological constraints while expanding available training examples. The balance between data-driven flexibility and knowledge-based constraint remains an active area of investigation, with different applications requiring different trade-offs between these complementary approaches.

\section{Current Limitations}
\label{sec:limitations}

Despite substantial progress in applying generative models to rare disease research, significant challenges remain before these approaches achieve routine clinical deployment. Validation frameworks lack standardization, privacy guarantees remain incomplete, computational requirements constrain accessibility, and the translation from research prototypes to regulatory-approved diagnostic tools faces substantial barriers. This section examines these limitations systematically to identify priorities for future development.

\subsection{Validation and Evaluation Challenges}

The absence of standardized validation protocols represents perhaps the most fundamental limitation constraining generative model deployment. Finetti et al. identified that over 40\% of reviewed publications provided insufficient quantitative detail regarding dataset characteristics, while fewer than 33\% reported both pre- and post-augmentation metrics \cite{finetti2025data}. This inconsistent reporting impedes meta-analysis, reproducibility assessment, and comparative evaluation across different generative approaches. Without consensus evaluation frameworks, researchers cannot reliably determine whether observed performance differences reflect genuine methodological advances or artifacts of dataset selection, hyperparameter choices, or evaluation protocols.

Current validation practices predominantly rely on downstream task performance—measuring classification accuracy, segmentation quality, or prediction metrics on models trained with synthetic data. While pragmatic, this approach provides only indirect evidence of synthetic data quality. Strong performance may result from the model memorizing training patterns rather than learning generalizable representations. Conversely, weak performance could indicate limitations in the downstream model architecture rather than synthetic data quality. Finetti et al. noted that biological plausibility assessment remains rare, with most studies evaluating synthetic samples solely through quantitative metrics rather than domain expert review \cite{finetti2025data}.

Visual inspection by clinical experts, employed successfully by Frid-Adar et al. for liver lesion validation \cite{fridadar2018gan}, does not scale to large synthetic datasets and introduces subjectivity that complicates systematic evaluation. Automated quality assessment metrics specific to biomedical domains remain underdeveloped. General-purpose image quality metrics such as Structural Similarity Index (SSIM) or Fréchet Inception Distance (FID) may not capture medically-relevant features that determine diagnostic utility. For electronic health records, statistical tests comparing marginal distributions fail to detect higher-order dependencies or clinically impossible code combinations.

The temporal dimension adds further complexity. For longitudinal health records or disease progression modeling, validation must assess not only cross-sectional realism but also the plausibility of temporal trajectories. A synthetic patient record might contain individually realistic diagnosis codes yet exhibit an impossible progression sequence—for example, recording diabetes complications before the initial diabetes diagnosis. Zhong et al. addressed this through conditioning mechanisms that respect temporal ordering \cite{zhong2024meddiffusion}, but systematic evaluation frameworks for longitudinal plausibility remain absent.

\subsection{Privacy and Security Vulnerabilities}

Researchers have promoted synthetic data generation as a privacy-preserving alternative to traditional de-identification approaches, yet substantial vulnerabilities persist. Mendes et al. documented that generative models remain susceptible to membership inference attacks, where adversaries determine whether the training set included a specific individual's data \cite{mendes2025synthetic}. Model inversion attacks can reconstruct approximations of training samples from model parameters, potentially exposing sensitive patient information. These attacks prove particularly concerning for rare diseases where small cohorts may enable reidentification even without direct attribute disclosure.

The distinction between fully synthetic and partially synthetic data critically impacts privacy guarantees. Fully synthetic datasets generated without direct reference to real individuals provide stronger privacy protection but may sacrifice analytical validity. Partially synthetic data, which combines real and fabricated elements, maintains higher fidelity to real-world distributions but introduces residual reidentification risks \cite{mendes2025synthetic}. The optimal balance depends on specific use cases, with different applications prioritizing privacy protection versus analytical accuracy differently.

Differential privacy offers theoretical guarantees by adding calibrated noise during training to bound the influence of any individual training sample. However, the privacy-utility trade-off proves challenging for rare diseases where small sample sizes make the impact of individual patients more pronounced. Strong privacy guarantees require noise levels that may substantially degrade model performance, while weak guarantees provide insufficient protection. Choi et al. demonstrated privacy-preserving EHR synthesis without explicit differential privacy mechanisms \cite{choi2017generating}, but the absence of formal privacy bounds limits confidence in the approach's robustness against sophisticated attacks.

Adversarial debiasing and fairness-aware algorithms represent emerging approaches to address representation biases in synthetic data. Mendes et al. noted that synthetic generation can either amplify or mitigate demographic biases present in training data \cite{mendes2025synthetic}. Without explicit fairness constraints, generative models may underrepresent minority populations or rare disease subtypes, perpetuating existing health disparities. Fairness evaluation frameworks for synthetic biomedical data remain nascent, with limited consensus on appropriate metrics or acceptable bias thresholds across different clinical contexts.

\subsection{Computational and Resource Barriers}

The computational demands of deep generative models create substantial accessibility barriers, particularly for smaller research institutions and clinical centers without dedicated machine learning infrastructure. GAN training typically requires GPU clusters, extensive hyperparameter tuning, and expertise in managing training instabilities such as mode collapse. Finetti et al. emphasized that these computational requirements often exceed the resources available to rare disease patient advocacy organizations and smaller medical centers that would most benefit from synthetic data expansion \cite{finetti2025data}.

Training times ranging from hours to days for individual models compound these challenges. Iterative refinement of architectures, hyperparameters, and conditioning mechanisms demands substantial computational budgets. Cloud computing infrastructure provides partial solutions but introduces data governance complexities when protected health information must be processed in third-party environments. The carbon footprint of large-scale model training also raises sustainability concerns as generative model deployment expands.

VAEs offer computational advantages relative to GANs, requiring fewer training iterations and exhibiting more stable convergence \cite{mendes2025synthetic}. However, this efficiency comes at the cost of reduced sample sharpness for imaging applications. Diffusion models, while achieving impressive generation quality, require multiple denoising steps during inference that slow sample generation compared to single-pass GAN or VAE synthesis \cite{zhong2024meddiffusion}. Recent advances in accelerated sampling partially address this limitation but do not eliminate the computational gap.

The expertise requirements extend beyond computational infrastructure to encompass deep learning engineering, domain knowledge integration, and clinical validation skills. Rare disease researchers with medical expertise may lack machine learning training, while machine learning researchers may lack the domain knowledge necessary to assess biological plausibility. This interdisciplinary challenge necessitates collaborative teams that prove difficult to assemble, particularly for ultra-rare conditions affecting dozens of patients globally.

\subsection{Clinical Translation and Regulatory Challenges}

The pathway from research prototypes to regulatory-approved clinical decision support tools remains poorly defined for synthetic data applications. Regulatory agencies including the U.S. Food and Drug Administration (FDA) and European Medicines Agency (EMA) lack established frameworks for evaluating medical devices trained on synthetic data. Key questions remain unresolved: What proportion of training data can be synthetic? How should synthetic data quality be assessed for regulatory purposes? What validation studies demonstrate that models trained on synthetic data generalize safely to real patients?

Clinical integration faces additional barriers beyond regulatory approval. Healthcare providers require interpretability and explainability to trust model recommendations, yet generative models function as black boxes with limited mechanistic transparency. Alsentzer et al. addressed this through attention mechanisms that identify which phenotypic features drive diagnostic predictions \cite{alsentzer2025few}, but such interpretability remains uncommon. Without clear explanations for why a model recommends specific diagnoses or treatments, clinicians may appropriately hesitate to incorporate these tools into practice.

Liability concerns compound adoption challenges. When a diagnostic error occurs with a model trained partially on synthetic data, determining responsibility proves complex. Is the synthetic data generator liable? The model developer? The clinician who applied the recommendation? These legal ambiguities discourage risk-averse healthcare institutions from deploying synthetic data-based systems despite potential clinical benefits.

The generalization of models across different clinical sites, patient populations, and temporal contexts requires extensive validation. Models trained at academic medical centers may fail when deployed at community hospitals with different patient demographics and disease prevalence. Alsentzer et al. demonstrated sustained performance across 12 clinical sites \cite{alsentzer2025few}, but such multi-site validation remains rare. Temporal validation, assessing whether models maintain performance as medical practice evolves and new treatments emerge, receives even less attention despite its critical importance for long-term deployment.

\subsection{Biological Plausibility and Domain Constraint Enforcement}

Generative models trained on purely data-driven objectives lack inherent mechanisms to enforce biological constraints, creating risks of generating physiologically impossible samples. A GAN might produce a medical image with tissue densities outside physically realizable ranges or generate an electronic health record with mutually exclusive diagnosis codes. While such samples may appear statistically consistent with training distributions, they violate domain knowledge that trained medical professionals recognize immediately.

Hybrid approaches that combine generative modeling with rule-based constraints offer partial solutions. Knowledge graphs encoding gene-disease-phenotype relationships, as employed by Alsentzer et al. \cite{alsentzer2025few}, guide generation toward biologically plausible combinations. Ontologies such as the Human Phenotype Ontology or Disease Ontology provide structured representations of valid medical concepts and their relationships, enabling post-generation filtering of invalid samples. However, comprehensively encoding all biological constraints proves infeasible given the complexity of human physiology and the incompleteness of current biomedical knowledge.

The challenge intensifies for novel or poorly characterized rare diseases where biological constraints remain incompletely understood. Generative models might produce samples that appear plausible based on limited existing data yet fail to capture emergent disease characteristics not represented in small training sets. This limitation proves particularly acute for the ultra-rare diseases where generative approaches promise the greatest potential impact but where biological understanding remains most limited.

Validation frameworks must distinguish between statistical fidelity and clinical utility. A synthetic dataset that perfectly reproduces training distribution statistics may nonetheless fail to capture subtle but clinically critical features. Conversely, synthetic data that exhibits some statistical discrepancies from real data might nonetheless provide sufficient signal for training clinically useful predictive models. Establishing which aspects of the real data distribution must be preserved and which can be approximated represents an ongoing challenge requiring close collaboration between machine learning researchers and domain experts.

\section{Conclusion and Future Directions}
\label{sec:conclusion}

Generative models have emerged as a critical tool for addressing data scarcity in rare disease research, enabling machine learning applications that would otherwise remain infeasible given limited patient cohorts. This review examined how Generative Adversarial Networks, Variational Autoencoders, and diffusion models adapt to the unique constraints of rare disease biomedical data across imaging, electronic health records, and genomic modalities. The evidence demonstrates quantifiable improvements in diagnostic performance, with sensitivity gains of 7.1 percentage points for liver lesion classification \cite{fridadar2018gan}, accuracy improvements to 85.9\% for brain MRI analysis \cite{mendes2025synthetic}, and PR-AUC increases of 6.83 percentage points for health risk prediction \cite{zhong2024meddiffusion}.

Despite these advances, substantial challenges temper optimism regarding immediate clinical deployment. Over 40\% of reviewed studies lack sufficient methodological detail for reproducibility assessment \cite{finetti2025data}, validation frameworks remain inconsistent across applications, and regulatory pathways for synthetic data-trained diagnostic tools remain undefined. Privacy vulnerabilities persist despite claims of synthetic data providing complete anonymization \cite{mendes2025synthetic}, while computational requirements create accessibility barriers for smaller research institutions. The biological plausibility of generated samples requires more rigorous evaluation than current performance-based validation provides.

\subsection{Synthesis of Key Findings}

The comparative analysis of generative model architectures reveals that optimal approaches depend critically on data modality, available computational resources, and dataset size. Medical imaging applications, which dominated 70.3\% of reviewed studies \cite{finetti2025data}, predominantly employ GANs due to their capacity for generating sharp, high-resolution images. The adversarial training objective effectively captures complex spatial dependencies in CT scans, MRI sequences, and histopathology slides, though at substantial computational cost and with risks of mode collapse. VAEs offer computational efficiency advantages for resource-constrained settings, trading some image sharpness for training stability. Diffusion models have demonstrated particular promise for sequential and temporal data, where maintaining consistency across time steps proves critical \cite{zhong2024meddiffusion}.

The integration of domain knowledge with data-driven generation represents a particularly important development for ultra-rare diseases. Alsentzer et al. demonstrated that knowledge graph-guided synthesis enables few-shot learning when fewer than 20 examples per disease are available \cite{alsentzer2025few}, achieving 40\% top-1 accuracy on real-world diagnostic tasks despite training primarily on 42,624 simulated patients. This paradigm shift from pure data-driven approaches to hybrid methods that incorporate structured biological knowledge may prove essential for the most data-constrained rare disease applications.

Dataset expansion rates ranging from threefold to tenfold have been consistently documented across modalities \cite{finetti2025data,mendes2025synthetic}, with the combination of classical and deep generative augmentation proving more effective than either approach alone \cite{fridadar2018gan}. However, the relationship between expansion magnitude and downstream performance remains nonlinear and application-dependent. Excessive reliance on synthetic data without sufficient real-world validation can lead to models that perform well on generated samples but fail to generalize to clinical deployment scenarios.

\subsection{Critical Gaps and Research Priorities}

Several critical gaps warrant prioritization in future research. First, standardized validation frameworks must be established to enable systematic comparison across generative approaches. Current inconsistencies in reporting dataset characteristics, expansion rates, and quality metrics impede progress toward identifying best practices. Community-driven initiatives to develop consensus evaluation protocols, similar to ImageNet for computer vision or MIMIC for clinical data, could accelerate methodological maturation.

Second, privacy-preserving techniques require theoretical foundations beyond empirical attack resistance. Differential privacy offers formal guarantees but introduces substantial utility degradation for small rare disease cohorts. Alternative approaches that provide provable privacy bounds while maintaining clinical utility represent high-priority research directions. The tension between strong privacy guarantees and the preservation of rare disease-specific characteristics that enable diagnostic model training demands careful theoretical and empirical investigation.

Third, interpretability and explainability mechanisms must advance beyond attention visualizations and feature importance scores. Regulatory approval and clinical adoption require that models provide mechanistic explanations for their predictions that clinicians can evaluate against domain knowledge. The integration of causal reasoning frameworks with generative modeling may enable models that not only generate realistic samples but also explain the biological mechanisms underlying their generations.

Fourth, multi-modal synthesis that simultaneously generates consistent data across imaging, genomic, and clinical record modalities remains largely unexplored. Real patients exhibit correlations between different data types—genetic variants influence phenotypic presentations that manifest in medical images and drive treatment decisions recorded in EHRs. Current approaches typically generate single modalities in isolation, missing opportunities to leverage cross-modal constraints that could improve sample quality and biological plausibility.

\subsection{Pathways to Clinical Translation}

The transition from research prototypes to clinically deployed diagnostic tools requires systematic addressing of regulatory, technical, and organizational challenges. Regulatory agencies must develop frameworks for evaluating synthetic data quality and determining acceptable proportions of training data that can be synthetic. Pilot studies demonstrating that regulatory-grade validation can be performed on synthetic data-trained models would establish precedents facilitating broader adoption.

Clinical validation studies must extend beyond single-institution evaluations to multi-site, multi-population assessments that demonstrate robust generalization. Alsentzer et al. showed sustained performance across 12 clinical sites \cite{alsentzer2025few}, providing a template for rigorous validation, yet such comprehensive evaluation remains rare. Prospective clinical trials comparing synthetic data-augmented models to standard-of-care diagnostic approaches would provide the evidence base necessary for confidence in clinical deployment.

Computational infrastructure development should prioritize accessibility for smaller research institutions and patient advocacy organizations. Cloud-based platforms that abstract away hardware requirements while maintaining data governance standards could democratize access to generative modeling capabilities. Pre-trained foundation models that can be fine-tuned on small disease-specific datasets might reduce computational demands compared to training from scratch.

Education and training initiatives must bridge the gap between medical domain expertise and machine learning technical skills. Rare disease researchers require accessible tools and workflows that do not presume deep learning expertise, while machine learning researchers need exposure to clinical contexts and biological constraints. Interdisciplinary training programs and collaborative infrastructure could accelerate progress by enabling effective teamwork across these disciplines.

\subsection{Future Technological Directions}

Several emerging technological directions promise to address current limitations. Federated learning approaches that enable model training across multiple institutions without centralized data pooling could expand effective sample sizes while respecting data governance constraints. Generative models trained through federated protocols could synthesize institution-specific data that respects local population characteristics while benefiting from multi-site knowledge aggregation.

Active learning frameworks that strategically select which samples to synthesize based on model uncertainty could improve sample efficiency. Rather than uniformly expanding all disease classes, these approaches would focus synthetic generation on underrepresented regions of feature space where additional data would most improve model performance. This targeted synthesis could achieve performance gains with fewer generated samples, reducing computational costs.

Continual learning mechanisms that enable models to update as new real-world data accumulates would address temporal validation challenges. Medical knowledge evolves as new treatments emerge and disease understanding advances; generative models must adapt rather than remaining static. Techniques that incorporate new information without catastrophically forgetting previous knowledge could enable long-term deployment that maintains relevance as medicine progresses.

The convergence of generative modeling with large language models and foundation models represents a particularly promising frontier. These models' broad knowledge of biomedical literature and structured medical data could inform generation processes, ensuring that synthetic samples respect known biological principles even when training data is extremely limited. The challenge lies in grounding generation in reliable domain knowledge rather than perpetuating errors or biases present in text corpora.

\subsection{Concluding Remarks}

Generative models have demonstrated clear value in expanding rare disease datasets and enabling machine learning applications that would otherwise fail due to insufficient training data. The quantitative performance improvements documented across multiple studies and clinical contexts provide compelling evidence that synthetic data generation represents a viable approach to data scarcity challenges. However, the path from research demonstrations to routine clinical deployment remains long, requiring coordinated efforts across technical development, regulatory framework establishment, clinical validation, and infrastructure development.

The rare disease community stands to benefit disproportionately from advances in generative modeling, given that data scarcity poses the most severe constraint on rare disease research relative to common conditions. Continued investment in these approaches, coupled with rigorous attention to validation, privacy, and clinical translation challenges, could materially accelerate progress toward more effective diagnostics and treatments for the millions of patients worldwide affected by rare diseases. The integration of synthetic data generation with traditional data collection and knowledge-guided learning frameworks offers a promising path forward for the field.


\bibliographystyle{IEEEtran}
\bibliography{refs}

\end{document}
